% Part: first-order-logic
% Chapter: models-theories
% Section: size-of-structures

\documentclass[../../../include/open-logic-section]{subfiles}

\begin{document}

\olfileid[pt]{fol}{mat}{siz}

\olsection{Expressando o Tamanho de \printtoken{P}{structure}}

\begin{explain}
Há algumas propriedades de estruturas que podemos expressar mesmo sem
usar os símbolos não lógicos de uma linguagem.  Por exemplo, há
!!{sentence}s que são verdadeiras em !!a{structure} se, e somente se, o
!!{domain} da !!{structure} tem pelo menos, no máximo ou exatamente um
certo número~$n$ de !!{element}s.
\end{explain}

\begin{prop}
A !!{sentence}
\begin{multline*}
  !A_{\ge n} \ident \lexists[x_1][\lexists[x_2][\dots\lexists[x_n][{}]]]\\
  \begin{aligned}
  (\eq/[x_1][x_2] \land {}
  \eq/[x_1][x_3] \land \eq/[x_1][x_4] \land \dots \land \eq/[x_1][x_n] \land {}\\
\eq/[x_2][x_3] \land \eq/[x_2][x_4] \land \dots \land {} \eq/[x_2][x_n] \land {} \\
\vdots\\
\eq/[x_{n-1}][x_n])
\end{aligned}
\end{multline*}
é verdadeira em !!a{structure}~$\Struct M$ se, e somente se, $\Domain M$
contém pelo menos $n$ !!{element}s. Consequentemente,
$\Sat{M}{\lnot !A_{\ge n+1}}$ se, e somente se, $\Domain M$ contém no
máximo~$n$ !!{element}s.

\end{prop}

\begin{prop}
A !!{sentence}
\begin{multline*}
!A_{= n} \ident \lexists[x_1][\lexists[x_2][\dots\lexists[x_n][{}]]] \\
\begin{aligned}
  (\eq/[x_1][x_2] \land {}
  \eq/[x_1][x_3] \land \eq/[x_1][x_4] \land \dots \land \eq/[x_1][x_n] \land {}\\
\eq/[x_2][x_3] \land \eq/[x_2][x_4] \land \dots \land {} \eq/[x_2][x_n] \land {} \\
\vdots\\
\eq/[x_{n-1}][x_n] \land {} \\
\lforall[y][(\eq[y][x_1] \lor \dots \lor \eq[y][x_n]]))
\end{aligned}
\end{multline*}
é verdadeira em !!a{structure}~$\Struct M$ se, e somente se, $\Domain M$
contém exatamente $n$ !!{element}s.
\end{prop}

\begin{prop}
Uma !!{structure} é infinita se, e somente se, é um modelo de
\[
\{!A_{\ge 1}, !A_{\ge 2}, !A_{\ge 3}, \dots \}.
\]
\end{prop}

Não há uma única sentença puramente lógica que seja verdadeira
em~$\Struct M$ se, e somente se, $\Domain M$ é infinito.  No entanto,
pode-se dar !!{sentence}s com !!{predicate}s não lógicos que só têm
modelos infinitos (embora nem toda !!{structure} infinita seja modelo
delas).  A propriedade de ser uma estrutura finita e a propriedade de
ser uma estrutura !!{nonenumerable} nem sequer podem ser expressas com
um conjunto infinito de !!{sentence}s.  Estes fatos seguem dos
teoremas da compacidade e de L\"owenheim--Skolem.

\end{document}
